\section{A shorter expansion of equality transport}
\label{sec:short-transport}

This appendix concerns the zero-normal scalar term language.  Write
\(\mathcal I(c,a;u.h)\) for the smart iterator, equal to \(a\) when
\(c\) is literally zero and otherwise equal to \(I(c,a;u.h)\).
Angle brackets denote capture-free substitution followed by zero
normalization.  No successor computation is performed by normalization.
The two arithmetic rules used here are
\[
 \frac{}{\Gamma\vdash v=v}\;{\mathsf{Refl}},
 \qquad
 \frac{\Gamma\vdash\mathcal I(c,a;u.h)=t
       \qquad\Gamma\vdash I(Sc,a;u.h)=r}
      {\Gamma\vdash r=h\langle t/u\rangle}\;{\mathsf{UE}}.
\]
All the following constructions work under an arbitrary retained context.
They use neither numerical induction nor a logical rule.

\paragraph{The inherited identity frame.}
For a term \(v\), put \(K(v)=I(S0,v;u.u)\).  The smart predecessor
iterate of \(K(v)\) is exactly \(v\).  Thus, given proofs
\(P:v=w\) and \(Q:v=r\), define
\[
 \mathsf{Root}(P,Q)
  :=\mathsf{UE}\bigl(P,\,
       \mathsf{UE}(Q,\mathsf{Refl}(K(v)))\bigr).
 \tag{R}
\]
The inner inference proves \(K(v)=r\), and the outer inference proves
\(r=w\).  This is the identity-frame construction from the incoming
\texttt{equality\_from\_iteration} package, reused here.  It adds exactly
two UE nodes and one reflexivity node, and uses each supplied proof once.

\paragraph{Congruence by two evaluations.}
Let \(h(u)\) be any fixed numerical template and let \(P:s=t\).
The template may contain any finite number of occurrences of \(u\),
including occurrences under further iterator binders.  Freshen the bound
variables to make every displayed substitution capture-free.  Put
\[
 H:=I(S0,s;u.h),\qquad
 h_s:=h\langle s/u\rangle,\qquad h_t:=h\langle t/u\rangle .
\]
The minor equation
\[
 e_s:=
 \mathsf{UE}\bigl(\mathsf{Refl}(s),\mathsf{Refl}(H)\bigr)
 \quad:\quad H=h_s
\]
is a one-step computation.  It already has exactly the left-hand side
needed by a second UE inference.  Therefore
\[
 \mathsf{Cong}_h(P)
   :=\mathsf{UE}(P,e_s)
   \quad:\quad h_s=h_t.
 \tag{C}
\]
Indeed, in the second inference the predecessor iterate is
\(\mathcal I(0,s;u.h)=s\); its major is \(s=t\), its minor is \(H=h_s\),
and its conclusion is \(h_s=h\langle t/u\rangle\).
This directly verifies the match, without an intervening equality
composition.

Consequently arbitrary numerical congruence adds exactly two UE nodes
and two reflexivity nodes to its supplied proof.  That proof occurs
once.  In particular, if \(P\) has \(p\) nodes, the displayed proof of
\(h_s=h_t\) has exactly \(p+4\) nodes.  The bound is independent of the
depth of \(h\) and the number of occurrences of \(u\), but not a bound
on serialized term or certificate size.

The incoming package instead first constructs both endpoint equations
\(H=h_s\) and \(H=h_t\), then composes them by (R).  Construction (C)
uses the first endpoint equation directly as the minor premise of UE.
The saving is in this explicit proof expansion; it does not constitute
a new axiom deletion or a claim of global proof-length optimality.

\paragraph{Crossed one-template transport.}
For \(P:s=t\) and \(Q:h_s=r\), define
\[
 \mathsf{Transport}_h(P,Q)
   :=\mathsf{Root}(\mathsf{Cong}_h(P),Q).
 \tag{T}
\]
Its conclusion is \(r=h_t\), which is precisely crossed one-template
transport.  Expanded without abbreviations, its proof is
\[
 \mathsf{UE}\!\left(
   \mathsf{UE}\!\left(
      P,\mathsf{UE}(\mathsf{Refl}(s),\mathsf{Refl}(H))
   \right),
   \mathsf{UE}\!\left(Q,\mathsf{Refl}(K(h_s))\right)
 \right).
\]
The inner right-hand inference proves \(K(h_s)=r\); the outer inference
then uses major \(h_s=h_t\) and the identity update to conclude
\(r=h_t\).

There are exactly four UE nodes and three reflexivity nodes in this
expansion.  Each supplied proof appears once.  Thus the complete tree
has \(p+q+7\) nodes when \(P,Q\) have \(p,q\) nodes.  The prior specified
expansion used six UE nodes and five reflexivity nodes, so this
construction removes four administrative inference nodes per replaced
transport instance.

\paragraph{Exact compilation count.}
Consider a proof tree in an intermediate calculus that already has all
the target rules and, in addition, the crossed one-template transport
rule.  Let the tree have \(N\) nodes, of which \(k\) are instances of
that extra rule.  Recursively replace every such instance by (T).
The resulting tree has exactly
\[
                         N+6k
\]
nodes.  Each source transport node is replaced by seven target nodes,
and neither of its recursively translated premise trees is duplicated.
This count is an external statement about finite proof trees.  If
other source rules are translated in a separate pass, their overhead
must be counted separately.  No uniform bound on certificate bytes,
checking time, or normal-form size follows from this node count.

\paragraph{Two-sided equation transport and sharing.}
For a supplied equation \(a(s)=b(s)\), applying (T) first to the template
\(a\) gives \(b(s)=a(t)\).  Applying (T) to the template \(b\) then gives
\(a(t)=b(t)\).  The second application needs a second occurrence of the
proof \(s=t\).  If the two supplied proof trees have \(p,q\) nodes, the
literal two-application expansion has
\[
                         2p+q+14
\]
nodes.  It is therefore incorrect to claim an additive constant
tree-size overhead for this particular two-sided expansion.  With
shared proof graphs the proof of \(s=t\) can be referenced twice
instead of copied, but that is a different representation and requires
its own checker and size convention.  The exact \(N+6k\) count above
applies to the one-template rule, for which no sharing assumption is
needed.

